\chapter*{Preface}
\addcontentsline{toc}{chapter}{Preface}
\markboth{Preface}{Preface}

\chapterstatus{working research monograph with completed synthetic numerical
verification and a proposed material program. No material validation is
claimed.}

This research monograph is a long-form working account of the
\texttt{ksdft2effmass} research program. It is not an institutionally
registered dissertation or a submitted manuscript. It combines accepted
synthetic software and numerical verification with proposed production and
material work. Its purpose is to develop one coherent narrative while
preserving the narrower scopes of articles, conference contributions, software
records, and computational evidence.

\section*{Research assistance and responsibility}
\index{artificial intelligence!research assistance}
\index{scientific rigor}


Large language models assisted with drafting, code development, analysis, and
editorial work. Their outputs were treated as provisional and evaluated under
the same source, verification, provenance, and evidentiary requirements as
other contributions. Scientific claims in this monograph rely on identified
literature and retained calculation or proof records rather than on the use of
an AI system. Responsibility for the selection, interpretation, and reporting
of the material remains with the author.

\clearpage
\section*{Audience}
\index{audience}
\index{pedagogy}


Although this research interest began with the development of a pedagogical
problem, I recognize the breadth of the resulting work. This manuscript serves
partly as a laboratory notebook and partly as a structured set of research
notes. It is intended to help advanced physics undergraduates and graduate
students understand the quantum-mechanical, solid-state, semiconductor,
computational, and mathematical concepts involved in the program.

The manuscript therefore includes a range of exercises developed to move a
reader from reading, through note-taking and computation, and ultimately toward
research. These exercises do not all carry the same evidentiary status. A
pedagogical derivation, controlled numerical experiment, proposed computational
workflow, and retained scientific calculation must remain distinguishable even
when they appear in one continuous narrative.

\section*{How to read the monograph}

The main narrative has four parts. Part~I asks what makes a reduced model
adequate and states the comparison and evidence rules used throughout the
book. Part~II develops the pristine bulk-silicon representation program, from
the first-principles parent through retained and localized representations to
bulk reduced-model selection. Part~III develops the phosphorus and boron
branches, aligned impurity extraction, lattice and continuum model classes,
and transferability questions. Part~IV collects the mathematical and proof
program needed to support the preceding scientific claims. The concluding
chapter then restates the dependency-ordered research sequence.

The appendices form a reference and analytical workbench. They contain detailed
notation, operator geometry, controlled finite-dimensional examples, reduction
route comparisons, completed one- and two-dimensional periodic benchmarks,
completed spin-space and one-dimensional defect benchmarks, a proposed
two-dimensional defect program, and a dedicated envelope-theory appendix
extracting the Luttinger--Kohn, Burt, and Ermoneit equations for the project.
They support the main argument while preserving the distinctions among
synthetic verification, literature analysis, proposed work, and material
validation.

\section*{Status and provenance}

The monograph explains the research program and reports evidence through
identified sources and retained records. Statements about future work are
proposals rather than results. A typeset formula is not by itself a proof, a
proof sketch is not a completed proof, a successful software test is not
scientific validation, and a generated figure is not a calculated result unless
its input, method, software environment, and provenance are identified. When a
calculation, source, or proof record changes the supported interpretation, the
narrative must be updated rather than treated as independent authority.

\section*{Objects kept distinct}

The discussion uses four related but noninterchangeable levels:

\begin{description}
  \item[Physical model.] The chosen description of the material system and its
    physical assumptions, including the Kohn--Sham parent problem and the
    substitutional-impurity branches.
  \item[Mathematical operator.] A map acting between identified state spaces,
    with a stated domain and codomain when these are material to the argument.
  \item[Finite representation.] A matrix obtained only after selecting finite
    subspaces, bases, ordering conventions, units, geometry, gauge, and energy
    reference.
  \item[Software artifact.] An implementation, serialized record, input file,
    workflow, or test that constructs or evaluates a representation.
\end{description}

Equality or subtraction at one level does not automatically define equality or
subtraction at another. In particular, pristine and doped matrices cannot be
subtracted meaningfully until their represented spaces and conventions have
been aligned.

\section*{Evidence vocabulary}
\index{evidence!vocabulary}


The monograph uses explicit status labels:

\begin{description}
  \item[Calculated result.] Produced by an identified calculation with retained
    provenance.
  \item[Literature value.] Taken from an identified and verified source.
  \item[Expected behavior.] A prediction or expectation that has not yet been
    observed in the project.
  \item[Illustrative example.] Explanatory material that is not research
    evidence.
  \item[Synthetic test data.] Constructed data used to verify software behavior,
    not output from a physical calculation.
  \item[Placeholder.] Material awaiting authoritative content or verification.
  \item[Proposed work.] Planned activity that has not been completed.
\end{description}

Chapter-level status statements apply to the prose that follows unless a more
specific label overrides them. Results chapters will identify their retained
sources rather than relying on narrative confidence or visual plausibility.

\section*{Error categories}
\index{error!categories}


Three error sources remain separate throughout the monograph. Parent-model
error concerns the difference between the selected physical parent and an
independent physical reference such as experiment or a higher-level theory.
Numerical or discretization error concerns approximation of the stated parent
mathematics by finite cutoffs, meshes, cells, stencils, and solver tolerances.
Model-reduction error concerns the loss incurred when replacing an aligned
parent representation by a reduced model class. These categories will not be
combined unless a later owning specification establishes compatible definitions
and a valid mathematical relationship among them.

\section*{Program scope}

The initial material system is silicon in the diamond structure. The declared
impurity branches are substitutional phosphorus and substitutional boron in
silicon. Quantum ESPRESSO remains responsible for electronic-structure
calculations, and Wannier90 remains responsible for Wannier localization. The
software in this repository begins from their retained outputs and addresses
representation, alignment, operator extraction, model reduction, evaluation,
and provenance.

The immediate bulk pilot and the later impurity program are related but not
identical scopes. Bulk development does not establish dopant results. Likewise,
a neutral periodic dopant supercell is not silently reinterpreted as an isolated
or ionized impurity. Each such interpretation requires its own declared branch
and evidence.

\citationtodo{medium priority}{Check the periodic-supercell versus isolated- or
ionized-impurity warning against a standard defect-physics source. Candidate:
Freysoldt et al. (2014), bibliography key \texttt{freysoldt2014} (already
present).}

\section*{Relationship to shorter publications}

Articles and presentations may select a narrower question, reorganize the
exposition, and adopt venue-specific terminology. Each publication must still
identify the evidence it uses and preserve the corresponding status,
provenance, and limitations. This monograph remains the broader working account
rather than evidence that any shorter publication has been reviewed or
accepted.
